\section{Results and Analysis}
\label{sec:results}

\subsection{Experimental Overview}
\label{sec:results:overview}

All results below are computed against the full official EDOS test split
\cite{kirk2023semeval}, $n=4{,}000$, at its real class distribution:
3{,}030 not-sexist examples (75.75\%) and 970 sexist examples (24.25\%).
We state this base rate up front because several of the accuracy figures
reported in this section are only interpretable against it: a
classifier that predicts ``not sexist'' unconditionally would already
score 0.758 accuracy by construction.

\subsection{Main Prompt-Design Grid}
\label{sec:results:grid}

Table~\ref{tab:main-grid} reports F1 for all 18 prompt variants across
the six evaluated models. Bold marks the best model for a given variant
(row); underline marks each model's own best-performing variant
(column).

\begin{table*}[t]
\centering
\footnotesize
\caption{F1 score by prompt variant (V1--V18) and model, full EDOS test
split ($n=4{,}000$). R = role, A = aspects, C = context, T =
chain-of-thought, F = few-shot ($k{=}2$ per class). Bold: best model per
row. Underline: best variant per model (column).
\textsuperscript{$\dagger$}llama3.1's V14 is omitted: 54.6\% of its
generations are safety refusals rather than classification attempts
(Sec.~\ref{sec:results:llama-refusal}). llama3.1's V8, V13, and V15
cells (F1 0.240, 0.361, 0.360) are reported but carry a related caveat:
39--49\% of their generations also fail to answer, mostly by the same
refusal mechanism (Table~\ref{tab:llama-cot}).}
\label{tab:main-grid}
\setlength{\tabcolsep}{4pt}
\begin{tabular}{l ccccc cccccc}
\toprule
& \multicolumn{5}{c}{Prompt Features} & \multicolumn{6}{c}{F1 Score} \\
\cmidrule(lr){2-6} \cmidrule(lr){7-12}
Variant & R & A & C & T & F & Mistral-7B & Phi-3-mini & Llama-3.1-8B & Qwen2.5-7B & Gemma-2-9B & Qwen3-8B \\
\midrule
V1 & -- & -- & -- & -- & -- & 0.467 & 0.138 & 0.441 & \textbf{0.537} & 0.448 & 0.481 \\
V2 & \checkmark & -- & -- & -- & -- & 0.454 & 0.256 & 0.472 & \textbf{0.535} & 0.444 & 0.483 \\
V3 & -- & \checkmark & -- & -- & -- & 0.490 & 0.413 & 0.550 & \textbf{0.570} & 0.501 & 0.543 \\
V4 & -- & -- & \checkmark & -- & -- & 0.453 & 0.056 & 0.466 & \textbf{0.493} & 0.454 & 0.477 \\
V5 & -- & -- & -- & \checkmark & -- & 0.496 & 0.480 & 0.412 & 0.509 & 0.443 & \textbf{0.553} \\
V6 & -- & -- & -- & -- & \checkmark & 0.567 & 0.517 & \underline{0.550} & \underline{0.585} & 0.530 & \textbf{0.590} \\
V7 & \checkmark & -- & -- & -- & \checkmark & \underline{0.567} & 0.533 & 0.519 & 0.583 & \underline{0.530} & \textbf{0.594} \\
V8 & -- & -- & -- & \checkmark & \checkmark & 0.471 & 0.438 & 0.240 & 0.571 & 0.503 & \underline{\textbf{0.627}} \\
V9 & \checkmark & \checkmark & -- & -- & \checkmark & 0.559 & \underline{\textbf{0.566}} & 0.471 & 0.559 & 0.514 & 0.562 \\
V10 & \checkmark & \checkmark & -- & \checkmark & -- & 0.461 & 0.478 & 0.497 & 0.543 & 0.454 & \textbf{0.544} \\
V11 & -- & \checkmark & \checkmark & \checkmark & -- & 0.483 & 0.466 & 0.508 & 0.556 & 0.471 & \textbf{0.559} \\
V12 & -- & \checkmark & -- & \checkmark & \checkmark & 0.462 & 0.461 & 0.501 & 0.568 & 0.496 & \textbf{0.602} \\
V13 & \checkmark & -- & -- & \checkmark & \checkmark & 0.473 & 0.448 & 0.361 & 0.576 & 0.482 & \textbf{0.614} \\
V14 & -- & -- & \checkmark & \checkmark & \checkmark & 0.447 & 0.461 & \textsuperscript{$\dagger$}-- & 0.571 & 0.500 & \textbf{0.618} \\
V15 & \checkmark & -- & \checkmark & \checkmark & \checkmark & 0.452 & 0.464 & 0.360 & 0.581 & 0.479 & \textbf{0.618} \\
V16 & \checkmark & \checkmark & \checkmark & \checkmark & -- & 0.474 & 0.473 & 0.495 & 0.537 & 0.454 & \textbf{0.546} \\
V17 & \checkmark & \checkmark & -- & \checkmark & \checkmark & 0.440 & 0.486 & 0.470 & 0.556 & 0.485 & \textbf{0.586} \\
V18 & \checkmark & \checkmark & \checkmark & \checkmark & \checkmark & 0.439 & 0.489 & 0.494 & 0.572 & 0.477 & \textbf{0.591} \\
\bottomrule
\end{tabular}
\end{table*}

Two patterns hold across all six models. First, the zero-shot baseline
(V1) is never any model's best variant, and the fullest five-component
prompt (V18) is not either: for every model, the top score comes from
a two- or three-component combination that includes few-shot examples.
Second, model behavior at V1 varies far more than the eventual ceiling
does: qwen2.5 opens at F1~0.537 with no prompt engineering at all, while
phi3 opens at F1~0.138, driven by a recall of just 0.082. It flags
almost nothing as sexist until given role framing, aspect decomposition,
or examples to work from. Qwen2.5 is the strongest model for every
V1--V4 (role/aspects/context-only) variant, but qwen3 becomes the
strongest model for every variant from V5 onward, once chain-of-thought
or few-shot enters the prompt. Qwen3's own best variant, V8
(chain-of-thought plus few-shot, F1~0.627), is also the single highest
score in the entire grid.

\subsection{Qualitative Finding: Safety-Alignment Interference in Llama-3.1-8B}
\label{sec:results:llama-refusal}

Llama-3.1-8B's V14 (context, chain-of-thought, few-shot; no role or
aspect framing) is omitted from Table~\ref{tab:main-grid} rather than
scored: 54.6\% of its 4{,}000 generations do not reach a parseable
YES/NO decision, reproduced identically across two independent runs
under this study's deterministic (greedy) decoding. Reading the raw
generations, these are not truncated or malformed outputs: they are
explicit safety refusals (e.g., ``I cannot create content that is
discriminatory or hateful''). Chain-of-thought is the trigger: no
variant without it shows a meaningful \texttt{fail\_ratio} for this
model. V14 is the worst single case of a pattern that runs across every
chain-of-thought variant llama3.1 was evaluated on, detailed below, not
an isolated one.

Scoring refusals as ``not sexist,'' the standard convention this study
uses elsewhere, V14 reaches 0.726 accuracy (almost exactly the test
split's 0.758 not-sexist base rate), with recall collapsing to 0.278
and F1 at 0.330. Restricted to the 45.4\% of rows the model actually
answers, F1 recovers to 0.525, in line with llama3.1's better-scoring
variants (V3 and V6, both F1~0.550). The two views tell different
stories: the model does not appear worse at the task when it engages,
but it declines to engage on over half the inputs under this specific
prompt structure. The answered subset is not a random sample either:
refusal plausibly correlates with how explicit the source text is. This
is a finding about llama3.1's safety alignment interacting with prompt
structure, not a gap in the evaluation; full derivation and the two-run
reproduction are documented alongside the raw refusal transcripts in the
project's supplementary material.

\begin{table}[t]
\centering
\footnotesize
\caption{\texttt{fail\_ratio} for llama3.1's ten other chain-of-thought
variants (none excluded from Table~\ref{tab:main-grid}; all stay below
the 0.5 exclusion gate). Refusal share is the fraction of a variant's
failed rows matching an explicit decline pattern (e.g.\ ``I cannot,''
``not able to''); the remainder are token-budget truncations, not
declines (see text).}
\label{tab:llama-cot}
\begin{tabular}{lcccc}
\toprule
Variant & Components & Aspects? & fail\_ratio & Refusal share \\
\midrule
V8  & T+F       & --         & 0.491 & 0.98 \\
V15 & R+C+T+F   & --         & 0.441 & 0.94 \\
V13 & R+T+F     & --         & 0.395 & 0.90 \\
V5  & T         & --         & 0.236 & 0.82 \\
V11 & A+C+T     & \checkmark & 0.214 & 0.31 \\
V12 & A+T+F     & \checkmark & 0.214 & 0.54 \\
V17 & R+A+T+F   & \checkmark & 0.171 & 0.63 \\
V16 & R+A+C+T   & \checkmark & 0.166 & 0.46 \\
V18 & R+A+C+T+F & \checkmark & 0.146 & 0.32 \\
V10 & R+A+T     & \checkmark & 0.121 & 0.45 \\
\bottomrule
\end{tabular}
\end{table}

Table~\ref{tab:llama-cot} reports \texttt{fail\_ratio} for llama3.1's
ten other chain-of-thought variants. Classifying each variant's failed
rows by whether they match an explicit refusal pattern reveals two
distinct causes, not one, and which dominates splits cleanly on whether
aspects is active. Every variant without aspects is refusal-dominant
(82--98\% of its failures are explicit declines, averaging around 100
characters each, matching V14's own transcripts). Every variant with
aspects is truncation-dominant or close to an even split (31--63\%
refusal); its non-refusal failures average 1{,}000--1{,}080 characters,
sitting right at this study's 220-token chain-of-thought budget
(Sec.~\ref{sec:experimental:generation}), and on inspection consist of
genuine step-by-step reasoning cut off mid-sentence before reaching a
decision, not a truncated refusal message. Aspects asks the model to
work through seven named linguistic markers before answering; for
llama3.1, that reasoning routinely does not finish inside a token budget
calibrated for the other five models' shorter prose reasoning. V8
(\texttt{fail\_ratio} 0.491), V13 (0.395), and V15 (0.441) are the three
most refusal-affected cells still reported in Table~\ref{tab:main-grid}
(F1 0.240, 0.361, and 0.360 respectively); their scores carry the same
caveat as V14's, just below rather than above the exclusion threshold: a
property of this specific model, worth knowing before comparing its
numbers to the other five.

\subsection{Statistical Significance of Prompt-Design Effects}
\label{sec:results:significance}

For each model we test every non-baseline variant against V1 with
exact McNemar's test (on the discordant-pair counts between the two
variants' predictions) and a paired bootstrap test on the F1 difference
(2{,}000 resamples), both Holm-corrected across the model's variants.
Table~\ref{tab:significance} summarizes the results.

\begin{table}[t]
\centering
\footnotesize
\caption{Significance of prompt-design effects vs.\ baseline V1,
Holm-corrected within each model ($\alpha=0.05$). ``Compared'' excludes
V1 itself and, for llama3.1, the omitted V14. ``McN.\ sig.''/``F1 sig.''
are the counts of variants significant under exact McNemar's test and
the paired F1-diff bootstrap, respectively (Sec.~\ref{sec:methodology:stats}).}
\label{tab:significance}
\setlength{\tabcolsep}{4pt}
\begin{tabular}{lccccc}
\toprule
Model & Compared & McN. sig. & F1 sig. & Best var. & Best $\Delta$F1 \\
\midrule
Mistral-7B & 17 & 17 & 12 & V7 & +0.100 \\
Phi-3-mini & 17 & 13 & 17 & V9 & +0.428 \\
Llama-3.1-8B & 16 & 16 & 16 & V6 & +0.110 \\
Qwen2.5-7B & 17 & 14 & 14 & V6 & +0.049 \\
Gemma-2-9B & 17 & 16 & 14 & V7 & +0.082 \\
Qwen3-8B & 17 & 16 & 15 & V8 & +0.146 \\
\midrule
Total & 101 & 92 & 88 & -- & -- \\
\bottomrule
\end{tabular}
\end{table}

88 of the 101 variant-vs-baseline comparisons across all six models
show a significant F1 difference, and 92 of 101 show a significant
McNemar effect: the large majority of this study's prompt-design choices
produce a real, statistically distinguishable effect on model behavior,
not sampling noise. The two tests disagree on a substantial fraction of
comparisons, in both directions (mistral: 17 McNemar-significant vs.\ 12
F1-diff-significant; phi3: 13 vs.\ 17). This divergence is expected,
since the two tests measure different quantities; phi3's V1$\to$V9 jump,
below, shows why.

Phi-3-mini's V1$\to$V9 jump is the clearest illustration. F1 rises from
0.138 to 0.566 (+0.428, by far the largest single effect in this study),
while accuracy moves only from 0.750 to 0.759 and the McNemar test does
not reach significance after correction ($p_{\text{holm}}=0.695$). The
reason is arithmetic, not a contradiction between tests: at V1, recall
is 0.082 and precision 0.421, implying on the order of 80 true positives
against 110 false positives out of 4{,}000 predictions. At V9, recall
reaches 0.647 and precision 0.502, implying roughly 628 true positives
against 622 false positives. The variant gains around 548 correct
sexist-flags and loses around 512 new false alarms, almost, but not
exactly, offsetting on a 4{,}000-row test set that is 76\% negative, so
raw accuracy (and the McNemar test built on it) barely moves. F1, which
is computed against the 970-row positive class rather than the full
test set, is not diluted by that majority-class denominator, so it
reflects the same underlying shift as a large, clearly significant
change. This is the concrete case for reporting both test families
rather than collapsing them into one number, and it is the opposite
pattern from every other model in this study, where McNemar catches
more significant comparisons than the F1-diff bootstrap does.

\subsection{Comparison Against Fine-Tuned and Classical Baselines}
\label{sec:results:baselines}

Table~\ref{tab:baselines} places each model's best training-free prompt
variant against a QLoRA-tuned version of the same model, a single
TF-IDF plus logistic-regression classifier (no LLM), and a
fine-tuned RoBERTa-base encoder.

\begin{table}[t]
\centering
\footnotesize
\caption{Best training-free prompting vs.\ fine-tuned and classical
baselines. Classical (TF-IDF + logistic regression): F1 0.617, accuracy
0.799. Encoder (RoBERTa-base): F1 0.735, accuracy 0.874. $\Delta$
columns are F1 relative to the classical baseline.}
\label{tab:baselines}
\begin{tabular}{lccc}
\toprule
Model & Best prompt F1 & QLoRA F1 & $\Delta$ vs.\ classical \\
\midrule
Mistral-7B & 0.567 (V7) & 0.788 & $-$0.050 \\
Phi-3-mini & 0.566 (V9) & 0.734 & $-$0.051 \\
Llama-3.1-8B & 0.550 (V6) & 0.795 & $-$0.067 \\
Qwen2.5-7B & 0.585 (V6) & 0.786 & $-$0.032 \\
Gemma-2-9B & 0.530 (V7) & 0.794 & $-$0.087 \\
Qwen3-8B & 0.627 (V8) & 0.779 & $+$0.010 \\
\bottomrule
\end{tabular}
\end{table}

The headline result of this comparison is not favorable to prompting.
\textbf{Five of the six models' best training-free prompt variant, after
an 18-point search over prompt design, scores below a bag-of-words
logistic regression classifier}: a baseline with no language model,
no GPU, and a fraction of a second of CPU inference per example. Only
qwen3 edges past it, by F1~+0.010, and even qwen3 loses to it on
accuracy (0.773 vs.\ 0.799). QLoRA fine-tuning and the RoBERTa encoder
are in a different tier entirely: every QLoRA-tuned model reaches
F1~0.734--0.795, and the encoder reaches F1~0.735, both comfortably
above the classical baseline and far above every prompted number.
Ranked by F1, the ordering is consistent across the whole study:
QLoRA~$\approx$~encoder~$\gg$~classical~$>$~training-free prompting,
with qwen3 as the one partial exception at the classical boundary. This
is the direct answer to this study's ``no non-LLM baseline'' gap: the
classical comparison reframes what the prompt-design results in
Sec.~\ref{sec:results:grid}--\ref{sec:results:significance} mean.
Prompt engineering produces large, statistically real differences
\emph{within} the space of training-free approaches, but that entire
space sits, for five of six models, below what a linear classifier over
word counts already achieves on this task.

\subsection{Error Analysis by Sexism Subtype}
\label{sec:results:subtype}

Table~\ref{tab:subtype} breaks down each model's best variant by the
four fine-grained sexism categories in the EDOS taxonomy
\cite{kirk2023semeval}, reporting the false-negative rate (share of
truly sexist examples in that category the model misses) for each.

\begin{table}[t]
\centering
\footnotesize
\caption{False-negative rate by sexism subtype, best variant per model.
1: threats/plans to harm/incitement ($n=89$). 2: derogation ($n=454$).
3: animosity ($n=333$). 4: prejudiced discussions ($n=94$).}
\label{tab:subtype}
\begin{tabular}{lcccccc}
\toprule
Model & Variant & 1 & 2 & 3 & 4 & Mean \\
\midrule
Mistral-7B & V7 & 0.011 & 0.057 & 0.147 & 0.181 & 0.099 \\
Phi-3-mini & V9 & 0.337 & 0.300 & 0.372 & 0.553 & 0.391 \\
Llama-3.1-8B & V6 & 0.079 & 0.064 & 0.129 & 0.255 & 0.132 \\
Qwen2.5-7B & V6 & 0.281 & 0.141 & 0.303 & 0.309 & 0.258 \\
Gemma-2-9B & V7 & 0.022 & 0.015 & 0.048 & 0.043 & 0.032 \\
Qwen3-8B & V8 & 0.303 & 0.150 & 0.237 & 0.372 & 0.266 \\
\bottomrule
\end{tabular}
\end{table}

Even at each model's own best variant, subtype-level performance
diverges far more than the aggregate F1 numbers in
Table~\ref{tab:main-grid} suggest. Gemma-2-9B misses under 5\% of
threats and derogation cases; phi3, at its own best variant, still
misses 55\% of ``prejudiced discussions'' and over a third of every
other category. Category 4 (prejudiced discussions, the smallest class
at $n=94$) is the hardest subtype for five of the six models. This
gap widens sharply when the worst rather than the best variant is
considered: phi3's worst variant (V4, context-only) misses 96--100\% of
sexist examples in every subtype, a near-total collapse into
predicting ``not sexist,'' while gemma2's worst variant (V8) never
exceeds a 16\% false-negative rate in any subtype. The spread between a
model's best-case and worst-case prompt is itself a per-model property,
not a constant: gemma2 is the most robust to a bad prompt choice in this
study, and phi3 the least.

\subsection{Sensitivity to Prompt-Component Order and Few-Shot Size}
\label{sec:results:sensitivity}

For each model's best main-grid variant, we re-evaluate under (a) an
alternate ordering of its prompt components, where the variant has more
than one component the harness can meaningfully reorder, and (b) an
alternate few-shot count ($k \in \{1,3,4\}$ per class, against the
main grid's default $k=2$). Table~\ref{tab:sensitivity} reports each as
an F1 delta from the main-grid value.

\begin{table}[t]
\centering
\footnotesize
\caption{Sensitivity of each model's best variant to component reordering
and few-shot count, as $\Delta$F1 from the Table~\ref{tab:main-grid}
value. n/a: variant not eligible for order sweep (see text).}
\label{tab:sensitivity}
\begin{tabular}{lcccc}
\toprule
Model (variant) & Order & $k{=}1$ & $k{=}3$ & $k{=}4$ \\
\midrule
Mistral-7B (V7) & n/a & $-$0.039 & +0.024 & +0.028 \\
Phi-3-mini (V9) & $-$0.032 & +0.024 & $-$0.019 & $-$0.092 \\
Llama-3.1-8B (V6) & n/a & +0.023 & $-$0.004 & +0.001 \\
Qwen2.5-7B (V6) & n/a & $-$0.031 & $-$0.005 & $-$0.002 \\
Gemma-2-9B (V7) & n/a & $-$0.018 & $-$0.011 & +0.011 \\
Qwen3-8B (V8) & +0.004 & $-$0.005 & $-$0.008 & +0.018 \\
\bottomrule
\end{tabular}
\end{table}

Order sensitivity applies only to phi3 (V9: role, aspects, few-shot) and
qwen3 (V8: chain-of-thought, few-shot). The other four models' best
variants either combine a single component with few-shot or combine two
components that render identically regardless of order under this
study's prompt template, so no meaningful reordering exists to test.
Where it does apply, the two models diverge sharply: qwen3 is stable
under reordering ($\Delta=+0.004$), while phi3 loses 0.032 F1 under the
alternate order, a real, if modest, position effect for this model
specifically.

The few-shot-size sweep carries a more general caveat. This study's
main-grid results all use $k=2$ examples per class, but that is not a
tuned optimum: four of six models score higher F1 at $k=4$ than at the
$k=2$ default (mistral +0.028, gemma2 +0.011, qwen3 +0.018, llama3.1
+0.001), meaning the headline numbers in Table~\ref{tab:main-grid}
understate what few-shot prompting can achieve for most of this study's
models. Phi-3-mini is the exception, and a stark one: F1 falls
monotonically as $k$ increases past the default, down to $-0.092$ at
$k=4$, driven by precision rising (0.530$\to$0.670) as recall collapses
(0.663$\to$0.366). Where the other five models either improve or stay
roughly flat with more examples, phi3 becomes systematically more
conservative: the same under-triggering behavior visible in its V1
baseline (Sec.~\ref{sec:results:grid}) re-emerges as few-shot count
grows, rather than being fixed by it.

\subsection{Sensitivity to Demonstration Selection}
\label{sec:results:seed-sensitivity}

The order and size sweeps above vary how the few-shot demonstrations are
arranged and how many are shown, but never which specific examples are
drawn: every main-grid variant samples its $k{=}2$-per-class
demonstrations once, with a fixed seed (42). Table~\ref{tab:seed-sensitivity}
reruns each model's best variant under three alternate seeds (1, 7, 123),
reporting the resulting F1 range as a delta from the main-grid value.

\begin{table}[t]
\centering
\footnotesize
\caption{Sensitivity of each model's best variant to which few-shot
demonstrations are drawn: F1 range across 3 alternate seeds (1, 7, 123),
as $\Delta$F1 from the Table~\ref{tab:main-grid} value (grid's own seed,
42).}
\label{tab:seed-sensitivity}
\begin{tabular}{lcc}
\toprule
Model (variant) & Grid F1 & $\Delta$F1 range \\
\midrule
Mistral-7B (V7) & 0.567 & $-$0.025 to +0.011 \\
Phi-3-mini (V9) & 0.566 & $-$0.235 to +0.007 \\
Llama-3.1-8B (V6) & 0.550 & $-$0.097 to $-$0.028 \\
Qwen2.5-7B (V6) & 0.585 & $-$0.037 to $-$0.011 \\
Gemma-2-9B (V7) & 0.530 & $-$0.034 to $-$0.012 \\
Qwen3-8B (V8) & 0.627 & $-$0.025 to $-$0.003 \\
\bottomrule
\end{tabular}
\end{table}

Four of six models stay within a narrow, consistently negative range
(gemma2, qwen2.5, and qwen3 within 0.02--0.04 F1; llama3.1 within
0.03--0.10 F1): every alternate seed tested scores at or below the main
grid's own draw. Phi-3-mini is again this study's outlier, and by the
widest margin observed in any sensitivity check reported here: F1 falls
from 0.566 to 0.330 under one alternate seed. This is not a parsing
failure, \texttt{fail\_ratio} is 0 on every seed tested, main grid
included, but a genuine behavioral shift: precision rises to 0.744 while
recall collapses to 0.212, the same conservative, under-triggering
pattern already observed as $k$ grows past the default
(Sec.~\ref{sec:results:sensitivity}) and in this model's V1 baseline
(Sec.~\ref{sec:results:grid}), now shown to also depend on which four
examples happen to be drawn. The main grid's own seed scores at or above
all three alternate seeds for four of six models (llama3.1, qwen2.5,
gemma2, qwen3); the two exceptions are mistral, where seeds 1 and 7 each
score marginally higher (+0.011, +0.002 F1), and phi3, where seed 7
scores 0.007 F1 higher despite this study's largest sensitivity-sweep
swing in the opposite direction coming from a different seed on the same
model. Seed 42 was fixed as an arbitrary default before this sweep was
run, not chosen for its result, but the pattern means
Table~\ref{tab:main-grid}'s headline numbers sit, for most models, closer
to the top than the middle of what a different demonstration draw would
produce.

\subsection{Computational Cost}
\label{sec:results:cost}

Table~\ref{tab:cost} reports wall-clock inference cost for the full
18-variant grid per model, recorded automatically as part of every run.

\begin{table}[t]
\centering
\footnotesize
\caption{Inference cost across the full 18-variant grid, per model.}
\label{tab:cost}
\begin{tabular}{lccc}
\toprule
Model & Total (h) & Mean (s/ex.) & Slowest variant \\
\midrule
Mistral-7B & 3.92 & 0.196 & V10, 25.9 min \\
Phi-3-mini & 3.19 & 0.160 & V12, 24.0 min \\
Llama-3.1-8B & 4.71 & 0.249 & V18, 31.5 min \\
Qwen2.5-7B & 1.51 & 0.075 & V11, 14.0 min \\
Gemma-2-9B & 6.26 & 0.313 & V18, 38.5 min \\
Qwen3-8B & 23.92 & 1.196 & V18, 146.8 min \\
\bottomrule
\end{tabular}
\end{table}

Qwen3-8B is an outlier by a wide margin: its full grid takes roughly
5--16$\times$ longer than any other model's, and its mean per-example
cost (1.196s) is 4--16$\times$ every other model's. This is a direct
consequence of its native reasoning mode, which this study leaves
enabled for chain-of-thought variants rather than capping its token
budget to match the other five models. Qwen3 is also this study's
best-performing model overall (Table~\ref{tab:main-grid},
\ref{tab:significance}); that performance is not free, and the
accuracy/cost trade-off it represents is a separate question from
whether it is the strongest model on accuracy alone.

\subsection{Confidence-Based AUC: A Second Comparison Against the Classical Baseline}
\label{sec:results:auc}

For the seven non-chain-of-thought variants (V1--V4, V6, V7, V9), the
model's decision is its first generated token, so a calibrated
P(sexist) score can be read directly off that token's YES/NO logits,
softmax-normalized between the two candidates, without generating
anything further. Chain-of-thought variants are excluded: their
decision only emerges after free-form reasoning, so no single generated
token stands in for it. We extract this score for all six models across
the full 4,000-example test split and score it against the true labels
with ROC-AUC. As a validity check, the confidence-implied call
(P(sexist) $\geq$ 0.5) agrees with the model's actual generated decision
on 94.5--99.7\% of rows for every model's variant in
Table~\ref{tab:auc}, confirming the first-token assumption holds on
real output, not only by construction.

\begin{table}[t]
\centering
\footnotesize
\caption{AUC from first-token P(sexist), best non-chain-of-thought
variant per model. Classical baseline (TF-IDF + logistic regression,
\texttt{predict\_proba}): AUC 0.848. $\Delta$ is AUC relative to the
classical baseline. \textsuperscript{$\dagger$}Qwen3-8B's true best
variant, V8 (F1 0.627, Table~\ref{tab:baselines}), is chain-of-thought
and has no single decision token this method can score; V7 is its best
scorable stand-in, not its true best variant.}
\label{tab:auc}
\begin{tabular}{lcccc}
\toprule
Model & Variant & F1 & AUC & $\Delta$ vs.\ classical \\
\midrule
Mistral-7B & V7 & 0.567 & 0.843 & $-$0.005 \\
Phi-3-mini & V9 & 0.566 & 0.810 & $-$0.038 \\
Llama-3.1-8B & V6 & 0.550 & 0.832 & $-$0.016 \\
Qwen2.5-7B & V6 & 0.585 & 0.812 & $-$0.036 \\
Gemma-2-9B & V7 & 0.530 & 0.875 & $+$0.027 \\
Qwen3-8B\textsuperscript{$\dagger$} & V7 & 0.594 & 0.838 & $-$0.010 \\
\bottomrule
\end{tabular}
\end{table}

Table~\ref{tab:auc} complicates the story in Table~\ref{tab:baselines}
rather than repeating it. On F1, qwen3 is the only model to beat the
classical baseline, and it does so with V8, a chain-of-thought variant
this AUC method cannot score at all. On AUC, qwen3's best scorable
variant loses to the classical baseline (0.838 vs.\ 0.848), while
gemma2, whose F1 (0.530) trails the classical baseline by the widest
margin in this study, is the only model whose AUC (0.875) beats it.
Which model, if any, ``beats'' the classical baseline depends on which
metric answers the question. Training-free prompting still fails to
consistently clear the classical baseline under either metric; the
exception to that just changes identity depending on which metric is
used.

\subsection{Confidence-Based Abstention: Trading Coverage for Precision}
\label{sec:results:abstention}

The same first-token P(sexist) score behind Table~\ref{tab:auc} also
supports a simple abstention policy: instead of forcing a decision on
every example, withhold a prediction whenever the score falls within a
symmetric band around the 0.5 decision boundary, deferring the closest
calls rather than answering them. For each model's Table~\ref{tab:auc}
variant, Table~\ref{tab:abstention} reports the resulting coverage (the
fraction of the test split still answered) and F1 on that covered subset
at a $\pm0.2$ band, against the same variant's full-coverage F1
(band$=$0) for reference.

\begin{table}[t]
\centering
\footnotesize
\caption{Confidence-based abstention at a $\pm0.2$ band around
P(sexist)$=$0.5, same variant per model as Table~\ref{tab:auc}. Coverage
is the fraction of the 4{,}000-example test split still answered after
abstention; F1 is computed on that covered subset only; $\Delta$F1 is
relative to the same variant's full-coverage (no-abstention) F1.}
\label{tab:abstention}
\begin{tabular}{lcccc}
\toprule
Model & Variant & Coverage & F1 & $\Delta$F1 \\
\midrule
Mistral-7B & V7 & 0.941 & 0.585 & $+$0.015 \\
Phi-3-mini & V9 & 0.628 & 0.691 & $+$0.109 \\
Llama-3.1-8B & V6 & 0.809 & 0.601 & $+$0.049 \\
Qwen2.5-7B & V6 & 0.971 & 0.593 & $+$0.007 \\
Gemma-2-9B & V7 & 0.959 & 0.541 & $+$0.011 \\
Qwen3-8B & V7 & 0.962 & 0.602 & $+$0.008 \\
\bottomrule
\end{tabular}
\end{table}

F1 improves for every model at this band. Beyond being calibrated in
aggregate (Table~\ref{tab:auc}), the confidence score tracks
per-example difficulty: the examples it flags as uncertain really are,
on average, more error-prone. The size of that improvement, and the
coverage it costs to get it, is not uniform. Phi-3-mini gains the most
($+$0.109 F1) but only by abstaining on 37\% of the test split (coverage
0.628), a policy this study's other five models never need to reach a
comparable gain: mistral, gemma2, qwen2.5, and qwen3 all retain over
94\% coverage at the same band while gaining under 0.02 F1. Llama-3.1-8B
falls between the two patterns ($+$0.049 F1 at 81\% coverage). This
mirrors a pattern already visible elsewhere in this study: phi3's errors
are not spread evenly across the test set but concentrated in a
recoverable subset (Sec.~\ref{sec:results:sensitivity},
Sec.~\ref{sec:results:seed-sensitivity}, and
Sec.~\ref{sec:results:subtype} all single out phi3 as an outlier among
the six models), and abstention is one more lens that finds the same
thing. As with Table~\ref{tab:auc}, this policy is evaluated only on the
seven non-chain-of-thought variants; whether the same abstention gain
holds for chain-of-thought variants, where no single token yields a
confidence score, is untested by this method.
